\documentclass[10pt,a4paper]{amsart}
\usepackage[margin=19mm]{geometry}
\usepackage{amsmath,amssymb,mathtools}
\usepackage[scr=rsfs,cal=euler]{mathalfa}
\usepackage{xfrac}
\usepackage[unicode,hidelinks,bookmarks=false]{hyperref}
\newcommand{\ie}{i.e.\ }
\newcommand{\cf}{cf.\ }
\newcommand{\resp}{respectively\ }
\newcommand{\Subsection}{\subsection}
\providecommand{\nobreakdash}{\nobreak\hbox{-}}
\setlength{\emergencystretch}{2em}
\multlinegap0pt
\usepackage{bm}
\usepackage{bbm}
\usepackage{dsfont}
\usepackage{xspace}
\usepackage{cancel}
\usepackage{mathtools}
\usepackage{amsthm}
\makeatletter
\@ifundefined{theorem}{\newtheorem{theorem}{Theorem}}{}
\@ifundefined{corollary}{\newtheorem{corollary}[theorem]{Corollary}}{}
\@ifundefined{lemma}{\newtheorem{lemma}[theorem]{Lemma}}{}
\makeatother
\makeatletter
\newcommand{\C}{{\mathbb{C}}}
\newcommand{\Ima}{\operatorname{Im }}
\newcommand{\N}{{\mathbb{N}}}
\newcommand{\R}{{\mathbb{R}}}
\newcommand{\Rea}{\operatorname{Re }}
\newcommand*{\defeq}{\mathrel{\vcenter{\baselineskip0.5ex \lineskiplimit0pt
 \hbox{\scriptsize.}\hbox{\scriptsize.}}}=}
\newcommand{\dist}{\operatorname{dist}}
\newcommand{\eqdef}{=\mathrel{\vcenter{\baselineskip0.5ex \lineskiplimit0pt
	\hbox{\scriptsize.}\hbox{\scriptsize.}}}}
\expandafter\ifx\csname subproof\endcsname\relax
\newenvironment{subproof}{\begin{proof}[Proof of claim]\renewcommand{\qedsymbol}{\subproofsymbol}}{%
               \end{proof}}
\fi
\newcommand{\uldelta}{\underline{\delta}}
\newcommand{\uleta}{\underline{\eta}}
\makeatother
\title[A finite-order answer to a problem of Erd\H{o}s on maximum m]{A finite-order answer to a problem of Erd\H{o}s on maximum modulus points}
\author{Leticia Pardo-Sim\'on, David J. Sixsmith}
\date{Source: arXiv:2607.09462v1 (CC BY 4.0)}
\begin{document}
\maketitle

\noindent\emph{Source and licence.} A finite-order answer to a problem of Erd\H{o}s on maximum modulus points. Version arXiv:2607.09462v1 (CC BY 4.0), distributed under the Creative Commons Attribution 4.0 International licence, as recorded in the arXiv licence metadata for this exact version and shown on the abstract page https://arxiv.org/abs/2607.09462v1. Full rights notice: licenses/arxiv-2607.09462-CC-BY-4.0.txt.\par\smallskip

\noindent\textit{Editorial scope.} Counted unit: the Lemma whose source label is \texttt{lem:narrow\_gate\_unique} in the article (source0.tex lines 660--664, source label \texttt{lem:narrow\_gate\_unique}), delivered together with the article's own complete original proof of it (source0.tex lines 665--1040, one \texttt{proof} environment of 376 lines). No mathematics is written, repaired or re-derived: statement and proof are the article's own text line for line. Required context retained uncounted: eq\_H (lines 218--221); eq:gG (lines 225--228); thm\_approx1.7 (lines 231--251); eq\_f (lines 237--244); eq\_S (lines 343--345); lem:psi\_basic\_estimates (lines 357--363); cor:psi\_large\_real (lines 388--406). Every definition, display and label of those lines is retained unchanged. Omitted from this package and not redistributed: 1--217, 222--224, 229--230, 245--342, 346--356, 364--387, 407--659, 1041--1912. Nothing inside a retained excerpt is omitted or altered: every retained body is delivered exactly as the article's lines are written, so no omission receipt of any kind accompanies this package. Citations: the retained excerpts carry no citation command at all, so no bibliography entry of the article is transcribed and no reference is redistributed. Reference adaptation: none was needed and none was made; every reference inside the delivered unit resolves inside it, so no unresolved reference or citation is produced. Printed slips kept literal: no printed word, symbol, hypothesis or number of the counted statement or of its proof was altered. Layout and class: the article's own class is \texttt{amsart} with packages including amsmath, amsfonts, amssymb, graphics, mathrsfs, amsthm, eurosym, verbatim, enumerate, subcaption, geometry, xcolor, enumitem, ulem; the wrapper is the lane's pinned \texttt{amsart} editorial wrapper at 10pt on A4, which restates the theorem-like environments the retained text needs and adds the article's own macro definitions that the retained text needs (\texttt{\textbackslash C}, \texttt{\textbackslash Ima}, \texttt{\textbackslash N}, \texttt{\textbackslash R}, \texttt{\textbackslash Rea}, \texttt{\textbackslash defeq}, \texttt{\textbackslash dist}, \texttt{\textbackslash eqdef}, \texttt{\textbackslash uldelta}, \texttt{\textbackslash uleta}), with extra wrapper packages (bm, bbm, dsfont, xspace, cancel, mathtools). The delivered numbering is the unit's own. Assets: no retained span contains a figure environment, a graphics-inclusion command, a PDF-page inclusion, a table or a TikZ picture; no image file is copied into this package, the package declares no asset dependency and no image file is redistributed with it. Licence: version arXiv:2607.09462v1 of this article is distributed under the Creative Commons Attribution 4.0 International licence, as recorded in the arXiv licence metadata for this exact version and shown on the abstract page https://arxiv.org/abs/2607.09462v1. A full rights notice is delivered at licenses/arxiv-2607.09462-CC-BY-4.0.txt. Characters: every retained excerpt is pure ASCII, so the delivered engine, XeLaTeX with its default Latin Modern text font, reports no missing character.
\begin{equation}
	\label{eq_H}
	H \defeq \{ x+iy \colon x> -14\log_+\vert y\vert \},
\end{equation}

\begin{equation}
	\label{eq:gG}
	g \circ \exp = \exp \circ G.
\end{equation} 

\begin{theorem}
	\label{thm_approx1.7}
	Let $(G, V, H)$ be a model, with $H$ fixed in \eqref{eq_H}, $-1\in V$, and assume that $G(-1)=1$.
	Let $g$ be as in \eqref{eq:gG}.
	%Let $(G, V, H)$ be a model, with $H$ fixed in \eqref{eq_H}, and let $g$ be as in \eqref{eq:gG}. 
	Then there exists $f\in \mathcal{B}$ such that
	\begin{equation}
		\label{eq_f}
		f(z) \defeq 
		\begin{cases}
			g(z)+h(z), &\text{for } z\in \exp(V), \\
			h(z),   &\text{otherwise.}
		\end{cases}
	\end{equation}
	Here $h\colon \C \rightarrow \C$ is such that, for some constant $M>0$ universal under this normalization, $\vert h(z)\vert \leq M/\vert z \vert $ for $|z| > 1$.

Moreover, suppose that \(T\defeq \exp(V)\) is symmetric with respect to the real axis and that, writing \(\Psi\colon T\to H\) for the map satisfying
	\(G=\Psi\circ\exp\), we have \(\Psi(T\cap\R)\subset\R\). Then \(f\) may be
	chosen so that \(f(\R)\subset\R\). Equivalently, since \(f\) is entire,
	\(f(\bar z)=\overline{f(z)}\) for all \(z\in\C\).
\end{theorem}

	\begin{equation}\label{eq_S}
		S \defeq \{ x + iy \in \C \colon |y| < 1/2 \},
	\end{equation}

\begin{lemma}\label{lem:psi_basic_estimates}
Let \(S\) be as in \eqref{eq_S}, and let \(\psi:S\to H\) be the conformal isomorphism with \(\psi(0)=1\) and \(\psi'(0)>0\). Then there are constants \(a_\psi,A_\psi>0\), depending only on \(H\), such that 
	\begin{equation}\label{eq:lem_aux1}
e^{a_\psi u}\psi(s)\leq\psi(s+u)\leq e^{A_\psi u}\psi(s) \quad \text{ for all } s\in\R \text{ and } u\geq0.
	\end{equation}
Moreover, for every \(s\in\R\), \(\max_{|y|<1/2}\Rea\psi(s+iy)=\psi(s)\).
\end{lemma}

	\begin{corollary}
	\label{cor:psi_large_real}Let \(S\) be as in \eqref{eq_S}, and let \(\psi:S\to H\) be the normalised conformal map from Lemma~\ref{lem:psi_basic_estimates}. For \(b\in\R\), set
	\[
	\Psi_b(\xi)\defeq \frac{\psi(b+\xi)}{\psi(b)},
	\qquad \xi\in S .
	\]
	Then, as \(b\to+\infty\),
	$
	\Psi_b(\xi)\longrightarrow e^{\pi \xi}
	$
	locally uniformly in \(S\), with convergence of derivatives. In particular,
	for every \(\rho\in(0,1/2)\),
	$
	\frac{\Rea \psi(b+it)}{\psi(b)}
	\longrightarrow
	\cos(\pi t)
	$
	in \(C^2([-1/2+\rho,1/2-\rho])\).
\end{corollary}

\begin{lemma}[Choice of the gate widths]\label{lem:narrow_gate_unique}
	There exists an admissible sequence \(\uleta=(\eta_n)_{n\in\N}\) such that, for every \(n\in\N\) and every \(\uldelta\in\Delta_{\uleta}\), the function \(z\mapsto |f_{\uleta,\uldelta}(e^z)|\) has a unique maximiser on \(\overline{I_{n,\uleta}(\uldelta)}\). This maximiser belongs to \(I_{n,\uleta}(\uldelta)\).
	
	Moreover, if there is \(K_0>0\) such that \(x_{n+1}-x_n\ge K_0\) for all \(n\), then, after the initial rescaling above, the sequence \(\uleta\) may be chosen to be constant: \(\eta_n=\eta_0\) for all \(n\), for some \(\eta_0>0\).
\end{lemma}

\begin{proof}
	We choose the numbers \(\eta_n\), and hence the widths of the gates, one at a time. Fix \(n\), and suppose that the previous
	gates have already been chosen. We first choose \(\eta_n\) small enough that
	\(\overline{D(w_n,20\eta_n)}\subset U\), that this disc is disjoint from all
	the other such discs already chosen, and that it meets no line
	\(\{\Rea z=x_m\}\), \(m\ne n\). We shall decrease \(\eta_n\) finitely many
	times below. For \(\uldelta\in\Delta_{\uleta}\), let
	$
	c_{n,\uldelta}
	=
	x_n+i(y_n-3\eta_n/2+\delta_n),
	$
	so that
	$
	I_{n,\uleta}(\uldelta)
	=
	\{c_{n,\uldelta}+i\eta_ns:|s|<1/2\}.
	$
	Set
	$
	\Omega_{n,\uldelta}
	=
	\eta_n^{-1}(V_{\uleta}(\uldelta)-c_{n,\uldelta}).
	$
	Let \(S=\{z:|\Ima z|<1/2\}\), let \(\psi:S\to H\) be the normalised map from
	Lemma~\ref{lem:psi_basic_estimates}, and write
	$
	\Phi_{\uleta,\uldelta}
	=
	\psi^{-1}\circ G_{\uleta,\uldelta}
	:
	V_{\uleta}(\uldelta)\to S .
	$
	Define
	\[
	F_{n,\uldelta}(\zeta)
	=
	\Phi_{\uleta,\uldelta}(c_{n,\uldelta}+\eta_n\zeta),
	\qquad
	\zeta\in\Omega_{n,\uldelta}.
	\]
	Thus \(F_{n,\uldelta}\) is the strip coordinate near the \(n\)-th gate,
	written in the rescaled variable. In this variable the gate is the fixed
	segment \(\{is:|s|<1/2\}\), and, as \(\eta_n\to0\), the rescaled domains
	\(\Omega_{n,\uldelta}\) converge in the Carath\'eodory kernel sense, with base point \(0\), to the two-sided slit plane
	\[
	\Omega
	=
	\C\setminus i\bigl((-\infty,-1/2]\cup[1/2,\infty)\bigr).
	\]
	Indeed, for the fixed \(n\), put
		\(d_n\defeq\min\{\dist(w_n,\partial U),\inf_{m\ne n}|x_m-x_n|\}>0\).
		Fix \(R>0\). Since \(|c_{n,\uldelta}-w_n|\le 3\eta_n/2\), the condition
		\((R+3/2)\eta_n<d_n\) implies, uniformly in
		\(\uldelta\in\Delta_{\uleta}\), that \(D(c_{n,\uldelta},R\eta_n)\subset U\)
		and that this disc meets no line \(\{\Rea z=x_m\}\), \(m\ne n\). Hence, after
		rescaling by \(z\mapsto (z-c_{n,\uldelta})/\eta_n\), the only removed part
		visible in \(D(0,R)\) is the \(n\)-th slit, and therefore
		\(\Omega_{n,\uldelta}\cap D(0,R)=\Omega\cap D(0,R)\) for all sufficiently
		small \(\eta_n\), uniformly in \(\uldelta\). Since \(R\) was arbitrary, the
		claimed Carath\'eodory kernel convergence follows.
	
	By Carath\'eodory kernel convergence for the two-sided slit, there are real
	numbers \(b_{n,\uldelta}\), tending to \(+\infty\) as \(\eta_n\to0\),
	uniformly for \(\uldelta\in\Delta_{\uleta}\), such that
	$
	F_{n,\uldelta}-b_{n,\uldelta}
	\longrightarrow
	\Theta
	$
	locally uniformly in \(\Omega\), with convergence of derivatives, where
	\[
	\Theta(\zeta)=
	\frac1\pi\operatorname{arsinh}(2\zeta)
	=
	\frac1\pi
	\log\left(2\left(\zeta+\sqrt{\zeta^2+1/4}\right)\right).
	\]
	Here \(\operatorname{arsinh}\) denotes the branch that is real on the real axis, the square-root branch is positive on the real axis, and the logarithm is real on the positive real axis. This map satisfies
	\[
	\Theta(0)=0,\qquad \Theta(\R)=\R,
	\qquad
	\Theta(x)\to\pm\infty
	\quad\text{as }x\to\pm\infty .
	\]
	
	On the gate of \(\Omega\), if \(\zeta=is\) and \(|s|<1/2\), then
	$
	\Theta(is)
	=
	\frac{i}{\pi}\arcsin(2s).
	$
	Consequently, as \(\eta_n\to0\), on compact subsegments of the gate,
	$
	F_{n,\uldelta}(is)
	=
	b_{n,\uldelta}
	+
	\frac{i}{\pi}\arcsin(2s)
	+ o(1),
	$
	uniformly in \(\uldelta\), with convergence of derivatives.
	
	Denote
	$
	\tau(s)\defeq \frac1\pi\arcsin(2s).
	$
	Since \(G_{\uleta,\uldelta}=\psi\circ\Phi_{\uleta,\uldelta}\), we have
	$
	\Rea G_{\uleta,\uldelta}(c_{n,\uldelta}+i\eta_ns)
	=
	\Rea \psi(F_{n,\uldelta}(is)).
	$
	By Corollary~\ref{cor:psi_large_real}, applied with
	\(b=b_{n,\uldelta}\), and by the preceding convergence of
	\(F_{n,\uldelta}-b_{n,\uldelta}\), we obtain
	\[
	P_{n,\uldelta}(s)
	\defeq
	\frac{
		\Rea G_{\uleta,\uldelta}(c_{n,\uldelta}+i\eta_ns)
	}{
		\psi(b_{n,\uldelta})
	}
	\longrightarrow
	\cos(\pi\tau(s))=
	\cos(\arcsin(2s))
	=
	2\sqrt{1/4-s^2}\eqdef	p(s)
	\]
	in \(C^2\) on compact subintervals of \((-1/2,1/2)\), uniformly in
	\(\uldelta\).

We also need a \(C^0\) version of this convergence at the two endpoints of
		the gate. 	We explain the upper endpoint; the lower one is analogous. In a
		neighbourhood of \(i/2\), use the coordinate
		\(\xi=\sqrt{-i(\zeta-i/2)}\), with the branch chosen so that the slit
		\(i[1/2,\infty)\) is opened into the boundary of a half-disc. In this
		coordinate the endpoint becomes a regular analytic boundary point, and, for
		small \(\eta_n\), the rescaled domains agree with the model domain near this
		point, uniformly in \(\uldelta\). Hence the boundary form of the
		Carath\'eodory convergence theorem gives uniform convergence of
		\(F_{n,\uldelta}-b_{n,\uldelta}\) to \(\Theta\) up to the endpoint. Applying
		the same argument at \(-i/2\), we obtain \(P_{n,\uldelta}\to p\) uniformly on
		the closed gate, where \(p(\pm1/2)=0\). Choose \(\rho>0\) so small that
	\(p(s)<1/4\) whenever \(1/2-\rho<|s|<1/2\). After decreasing \(\eta_n\),
	we therefore have
	\[
	P_{n,\uldelta}(s)<\frac12
	\qquad
	\text{whenever }1/2-\rho<|s|<1/2,
	\]
	uniformly in \(\uldelta\). It remains to look at the compact interval
	\[
	K_\rho\defeq[-1/2+\rho,1/2-\rho].
	\]
	There the convergence \(P_{n,\uldelta}\to p\) is in \(C^2\). Since \(p\) has
	a unique non-degenerate maximum at \(0\), the functions
	\(P_{n,\uldelta}\) also have exactly one maximum on \(K_\rho\), after
	decreasing \(\eta_n\). Moreover, since \(p(0)=1\), this maximum has value
	larger than \(3/4\), after decreasing \(\eta_n\). The endpoint estimate
	therefore rules out maxima in the two endpoint regions. Hence
	\(P_{n,\uldelta}\), and therefore
	$
	s\mapsto
	\Rea G_{\uleta,\uldelta}(c_{n,\uldelta}+i\eta_ns),
	$
	has a unique maximum on \((-1/2,1/2)\). Since \(r\mapsto e^r\) is strictly
	increasing, the same is true for
	$
	s\mapsto
	|g_{\uleta,\uldelta}(e^{c_{n,\uldelta}+i\eta_ns})|.
	$
	
	To conclude the proof, we must pass from the model function to the entire
	function. By Theorem~\ref{thm_approx1.7}, applied in logarithmic
	coordinates,
	$
	|h_{\uleta,\uldelta}(e^z)|\le Me^{-\Rea z}
	$
	whenever \(\Rea z>L\). Since the gate lies on the line \(\Rea z=x_n\), and
	since \(x_n-1>L\), this gives in particular
	that $
	|h_{\uleta,\uldelta}(e^{c_{n,\uldelta}+i\eta_ns})|
	\le Me^{-x_n},$ whenever $-1/2\le s\le1/2$. Let
	\[
	u(s)=\Rea G_{\uleta,\uldelta}(c_{n,\uldelta}+i\eta_ns).
	\]
	The estimates used below are consequences of the preceding
		\(C^2\)-convergence \(P_{n,\uldelta}\to p\), together with the definition of
		\(P_{n,\uldelta}\). Choose \(0<\alpha<1/2-\rho\) and
		\(d,\kappa,c>0\), independent of \(\uldelta\), such that, for the limit
		function \(p\), we have \(p''\le -2\kappa\) and \(p\ge 2c\) on
		\([-\alpha,\alpha]\), while \(p\le1-3d\) on
		\(K_\rho\setminus(-\alpha,\alpha)\). Combining this with the endpoint
		estimate above, and decreasing \(\eta_n\) if necessary, we obtain uniformly
		in \(\uldelta\) that the unique maximum point \(s_*=s_*(\uldelta)\) lies in
		\((-\alpha,\alpha)\), that
		\(P_{n,\uldelta}(s)\le P_{n,\uldelta}(s_*)-d\) for
		\(\alpha\le |s|<1/2\), that \(P_{n,\uldelta}''(s)\le-\kappa\) for
		\(|s|\le\alpha\), and that \(P_{n,\uldelta}(s)\ge c\) for
		\(|s|\le\alpha\). Since \(u\) is \(\psi(b_{n,\uldelta})\) times
		\(P_{n,\uldelta}\) by the definition of \(P_{n,\uldelta}\), these give
		\[
			\begin{aligned}
				u(s)&\le u(s_*)-d\psi(b_{n,\uldelta})
				&&\text{for } \alpha\le |s|<1/2,\\
				u''(s)&\le -\kappa\psi(b_{n,\uldelta})
				&&\text{for } |s|\le\alpha,\\
				u(s)&\ge c\psi(b_{n,\uldelta})
				&&\text{for } |s|\le\alpha .
		\end{aligned}
		\]
	Let
	\(r_\alpha\defeq(1/2-\alpha)/2\). After decreasing $\eta_n$, the points $c_{n,\uldelta}+i\eta_ns$, with $\operatorname{dist}(s,[-\alpha,\alpha])\le r_\alpha$, belong to $V_{\uleta}(\uldelta)$ and have real part at least $x_{n-1}$, uniformly in $\uldelta$. Hence Cauchy's estimates applied to the holomorphic function
	\(s\mapsto h_{\uleta,\uldelta}(e^{c_{n,\uldelta}+i\eta_ns})\) give, for
	\(k=0,1,2\),
	\[
	\left|
	\frac{d^k}{ds^k}
	h_{\uleta,\uldelta}
	\bigl(e^{c_{n,\uldelta}+i\eta_ns}\bigr)
	\right|
	\le
	C_\alpha Me^{-x_n},
	\qquad |s|\le\alpha,
	\]
	with \(C_\alpha\) independent of \(\uldelta\). 
	On \([-\alpha,\alpha]\), we have
	\[
	\begin{aligned}
		f_{\uleta,\uldelta}
		\bigl(e^{c_{n,\uldelta}+i\eta_ns}\bigr)
		&=
		e^{G_{\uleta,\uldelta}(c_{n,\uldelta}+i\eta_ns)}
		\left(
		1+
		h_{\uleta,\uldelta}
		\bigl(e^{c_{n,\uldelta}+i\eta_ns}\bigr)
		e^{-G_{\uleta,\uldelta}(c_{n,\uldelta}+i\eta_ns)}
		\right).
	\end{aligned}
	\]
	The term in parentheses is a \(C^2\)-small perturbation of \(1\), uniformly
	in \(\uldelta\), as \(\eta_n\to0\). Indeed, the \(h\)-term and its first two
	\(s\)-derivatives are \(O(Me^{-x_n})\), while on \([-\alpha,\alpha]\)
	$
	\Rea G_{\uleta,\uldelta}(c_{n,\uldelta}+i\eta_ns)
	\ge
	c\psi(b_{n,\uldelta})
	$
	and
	\[
	\frac{d^j}{ds^j}
	G_{\uleta,\uldelta}(c_{n,\uldelta}+i\eta_ns)
	=
	O(\psi(b_{n,\uldelta}))
	\qquad (j=1,2).
	\]
	Thus differentiating the factor
	\(e^{-G_{\uleta,\uldelta}(c_{n,\uldelta}+i\eta_ns)}\) produces only polynomial
	factors in \(\psi(b_{n,\uldelta})\), whereas
	$
	e^{-\Rea G_{\uleta,\uldelta}(c_{n,\uldelta}+i\eta_ns)}
	\le
	e^{-c\psi(b_{n,\uldelta})}.
	$
	Since \(\psi(b_{n,\uldelta})\to+\infty\), the perturbation tends to \(0\) in
	\(C^2([-\alpha,\alpha])\), uniformly in \(\uldelta\). Therefore
	$
	s\mapsto
	\log\left|
	f_{\uleta,\uldelta}
	\bigl(e^{c_{n,\uldelta}+i\eta_ns}\bigr)
	\right|
	$
	is a \(C^2\)-small perturbation of \(u(s)\) on \([-\alpha,\alpha]\). Since
	\(u''\le-\kappa\psi(b_{n,\uldelta})\) there, this function is strictly
	concave after decreasing \(\eta_n\). Hence
	$
	s\mapsto
	\left|
	f_{\uleta,\uldelta}
	\bigl(e^{c_{n,\uldelta}+i\eta_ns}\bigr)
	\right|
	$
	has at most one maximum in \((-\alpha,\alpha)\).
	
	It remains to rule out points outside \((-\alpha,\alpha)\). For
	\(\alpha\le |s|<1/2\),
	\[
	\left|
	e^{G_{\uleta,\uldelta}(c_{n,\uldelta}+i\eta_ns)}
	\right|
	\le
	e^{-d\psi(b_{n,\uldelta})}
	\left|
	e^{G_{\uleta,\uldelta}(c_{n,\uldelta}+i\eta_ns_*)}
	\right|.
	\]
	Moreover
	\[
	\left|
	e^{G_{\uleta,\uldelta}(c_{n,\uldelta}+i\eta_ns_*)}
	\right|
	=
	e^{u(s_*)}
	\ge
	e^{c\psi(b_{n,\uldelta})}.
	\]
	Thus, since the additive error coming from \(h_{\uleta,\uldelta}\) is bounded
	by \(CMe^{-x_n}\), after decreasing \(\eta_n\) no point with
	\(\alpha\le |s|<1/2\) can give a value as large as the value at \(s_*\).
	Consequently the maximum on the open gate is attained in
	\((-\alpha,\alpha)\), where it is unique.
	
	Finally, let \(z\) be one of the two endpoints of
	\(\overline{I_{n,\uleta}(\uldelta)}\). Then
	\(z\in L_{n,\uleta}(\delta_n)\), and hence
	\(e^z\notin\exp(V_{\uleta}(\uldelta))\). By the definition
	\eqref{eq_f} of the approximating function, the model term is absent at
	\(e^z\), and therefore
	\[
	|f_{\uleta,\uldelta}(e^z)|
	=
	|h_{\uleta,\uldelta}(e^z)|
	\le
	Me^{-x_n}.
	\]
	At the interior point corresponding to \(s_*\), however,
	\[
	\left|
	f_{\uleta,\uldelta}
	\bigl(e^{c_{n,\uldelta}+i\eta_ns_*}\bigr)
	\right|
	\ge
	e^{c\psi(b_{n,\uldelta})}
	-
	CMe^{-x_n}.
	\]
	After decreasing \(\eta_n\), this is larger than \(Me^{-x_n}\). Therefore
	neither endpoint can be a maximiser. Since the open gate has a unique
	maximiser, the closed gate
	\(\overline{I_{n,\uleta}(\uldelta)}\) also has a unique maximiser, and this
	maximiser lies in \(I_{n,\uleta}(\uldelta)\).
	
	Proceeding inductively over \(n\), and decreasing each \(\eta_n\) as above,
	gives the required admissible sequence.
	
	It remains to justify the final assertion. Suppose that \(x_{n+1}-x_n\ge K_0>0\) for all \(n\). Since \(y_n\in(-\pi,\pi]\), the slopes of the line segments forming \(\gamma\) are bounded in terms of \(K_0\). After the initial rescaling, the same is true, with a uniform bound, for the two sides of the initial trapezoid \(T\). Hence there is \(r_0>0\) such that \(D(w_n,r_0)\subset U\) for every \(n\). Choose \(\eta_0>0\) so small that $0<\eta_0<1/10$, \(20\eta_0<r_0\) and \(40\eta_0<K_0\). Then the constant sequence \(\eta_n=\eta_0\) is admissible.
	
	We claim that, after decreasing \(\eta_0\) if necessary, the uniqueness
	argument above holds uniformly for all \(n\) and all
	\(\uldelta\in\prod_n[0,3\eta_0]\). Fix \(R>0\). Since
	\(|c_{n,\uldelta}-w_n|\le3\eta_0/2\), if \(\eta_0\) is sufficiently small,
	depending only on \(R,r_0,K_0\), then
	\(D(c_{n,\uldelta},R\eta_0)\subset U\) and this disc meets no line
	\(\{\Rea z=x_m\}\) with \(m\ne n\), for every \(n\) and every
	\(\uldelta\). Hence, in the rescaled coordinate
	\(\zeta=(z-c_{n,\uldelta})/\eta_0\), the only boundary pieces visible in
	\(D(0,R)\) are the two parts of the \(n\)-th slit. Thus
	\[
	\eta_0^{-1}\bigl(V_{\uleta}(\uldelta)-c_{n,\uldelta}\bigr)
	\longrightarrow
	\Omega=\C\setminus i\bigl((-\infty,-1/2]\cup[1/2,\infty)\bigr)
	\]
	in the Carath\'eodory kernel sense, uniformly in \(n\) and
	\(\uldelta\). The corresponding convergence of
	\(F_{n,\uldelta}-b_{n,\uldelta}\) to \(\Theta\), including convergence of
	derivatives on compact subsets of \(\Omega\), is therefore also uniform in
	\(n\) and \(\uldelta\). Moreover \(b_{n,\uldelta}\to+\infty\) uniformly as
	\(\eta_0\to0\), by the same argument as in the preceding single-gate case. Since \(x_n\ge x_1>L+3\), the approximation error
	\(Me^{-x_n}\) is uniformly bounded by \(Me^{-x_1}\). Hence all estimates in
	the uniqueness proof hold simultaneously for every gate once \(\eta_0\) is
	sufficiently small. This proves the constant-width assertion.
\end{proof}

\end{document}
